% problem-000053 5 同位素丰度与Lu-Hf定年
\section*{5. 同位素丰度与Lu-Hf定年}
\textit{下列为分问。}

\subsection*{5-1}
已知 Cl 有两种同位素 35Cl 和 37Cl，⼆者丰度⽐为 0.75:0.25；Rb 有 85Rb 和 87Rb，⼆者丰度⽐为0.72:0.28。

\subsection*{5-1}
-1 写出⽓态中同位素组成不同的RbCl分⼦。

\subsection*{5-1}
-2 这些分⼦有⼏种质量数？写出质量数，并给出其⽐例。

\subsection*{5-2}
年代测定是地质学的⼀项重要⼯作。Lu-Hf法是上世纪80年代随着等离⼦发射光谱、质⼦谱等技术发展⽽建⽴的⼀种新断代法。Lu有两种天然同位素：176Lu和 $^ { 1 7 7 } \mathrm { L u } _ { \mathrm { o } }$ Hf有六种天然同位素：176Hf、177Hf、178Hf、179Hf、180Hf 和 181Hf。176Lu 发⽣ $\lvert \beta \rvert$ 衰变⽣成176Hf，半衰期为 $3 . 7 1 6 \times 1 0 ^ { 1 0 }$ 年。177Hf为稳定同位素且⽆放射性来源。地质⼯作者获得⼀块岩⽯样品，从该样品的不同部位取得多个样本进⾏分析。其中的两组有效数据如下：样本 1：176Hf 与 177Hf 的⽐值为 0.28630（原⼦⽐，记为 176Hf/177Hf），176Lu/177Hf 为 0.42850。

样本 2：176Hf/177Hf 为 0.28239，176Lu/177Hf 为 0.01470。

（⼀级反应，物种含量�随时间�变化的关系式： $\begin{array} { r } { c = c _ { 0 } e ^ { - k t } \underline { { \mathbb { H } } } \underline { { \hat { \mathbf { \Pi } } } } | \mathrm { l n } \frac { c } { c _ { 0 } } = - k t } \end{array}$ ，其中 $\iota _ { c _ { 0 } }$ 为起始含量）

\subsection*{5-2}
-1 写出 $^ \mathrm { 1 7 6 } \mathrm { L u }$ 发⽣�衰变的核反应⽅程式（标出核电荷数和质量数）。

\subsection*{5-2}
-2 计算 $^ { 1 7 6 } \mathrm { L u }$ 衰变反应速率常数�。

\subsection*{5-2}
-3 计算该岩⽯的年龄。

\subsection*{5-2}
-4 计算该岩⽯⽣成时 176Hf/177Hf的⽐值。
